\documentclass[10pt,a4paper]{amsart}
\usepackage[margin=19mm]{geometry}
\usepackage{amsmath,amssymb,mathtools}
\usepackage[scr=rsfs,cal=euler]{mathalfa}
\usepackage{xfrac}
\usepackage[unicode,hidelinks,bookmarks=false]{hyperref}
\newcommand{\ie}{i.e.\ }
\newcommand{\cf}{cf.\ }
\newcommand{\resp}{respectively\ }
\newcommand{\Subsection}{\subsection}
\providecommand{\nobreakdash}{\nobreak\hbox{-}}
\setlength{\emergencystretch}{2em}
\multlinegap0pt
\usepackage{bm}
\usepackage{bbm}
\usepackage{dsfont}
\usepackage{xspace}
\usepackage{cancel}
\usepackage{mathtools}
\usepackage{cleveref}
\usepackage{amsthm}
\makeatletter
\@ifundefined{lemma}{\newtheorem{lemma}{Lemma}}{}
\makeatother
\title[Pullback exponential attractors for 3D non-autonomous Boussi]{Pullback exponential attractors for 3D non-autonomous Boussinesq equations with fractional Laplacian}
\author{Tong Kou, Hui Liu, Xiuqing Wang, Jie Xin}
\date{Source: arXiv:2609.18114v1 (CC BY 4.0)}
\begin{document}
\maketitle

\noindent\emph{Source and licence.} Pullback exponential attractors for 3D non-autonomous Boussinesq equations with fractional Laplacian. Version arXiv:2609.18114v1 (CC BY 4.0), distributed under the Creative Commons Attribution 4.0 International licence, as recorded in the arXiv licence metadata for this exact version and shown on the abstract page https://arxiv.org/abs/2609.18114v1. Full rights notice: licenses/arxiv-2609.18114-CC-BY-4.0.txt.\par\smallskip

\noindent\textit{Editorial scope.} Counted unit: the Lemma whose source label is \texttt{lem0301} in the article (source0.tex lines 312--322, source label \texttt{lem0301}), delivered together with the article's own complete original proof of it (source0.tex lines 323--665, one \texttt{proof} environment of 343 lines). No mathematics is written, repaired or re-derived: statement and proof are the article's own text line for line. Required context retained uncounted: eq0101 (lines 87--94); ass:forcing (lines 153--158). Every definition, display and label of those lines is retained unchanged. Omitted from this package and not redistributed: 1--86, 95--152, 159--311, 666--1992. Nothing inside a retained excerpt is omitted or altered: every retained body is delivered exactly as the article's lines are written, so no omission receipt of any kind accompanies this package. Citations: the retained excerpts carry no citation command at all, so no bibliography entry of the article is transcribed and no reference is redistributed. Reference adaptation: none was needed and none was made; every reference inside the delivered unit resolves inside it, so no unresolved reference or citation is produced. Printed slips kept literal: no printed word, symbol, hypothesis or number of the counted statement or of its proof was altered. Layout and class: the article's own class is \texttt{elsarticle} with packages including amssymb, amsthm, mathrsfs, amsmath, lineno, geometry, xcolor, hyperref, cleveref; the wrapper is the lane's pinned \texttt{amsart} editorial wrapper at 10pt on A4, which restates the theorem-like environments the retained text needs and adds the article's own macro definitions that the retained text needs (none), with extra wrapper packages (bm, bbm, dsfont, xspace, cancel, mathtools, cleveref). The delivered numbering is the unit's own. Assets: no retained span contains a figure environment, a graphics-inclusion command, a PDF-page inclusion, a table or a TikZ picture; no image file is copied into this package, the package declares no asset dependency and no image file is redistributed with it. Licence: version arXiv:2609.18114v1 of this article is distributed under the Creative Commons Attribution 4.0 International licence, as recorded in the arXiv licence metadata for this exact version and shown on the abstract page https://arxiv.org/abs/2609.18114v1. A full rights notice is delivered at licenses/arxiv-2609.18114-CC-BY-4.0.txt. Characters: every retained excerpt is pure ASCII, so the delivered engine, XeLaTeX with its default Latin Modern text font, reports no missing character.
\begin{equation}\label{eq0101}
\begin{cases}
\partial_{t}u+\nu(-\Delta)^{\alpha}u+(u\cdot\nabla)u+\nabla p = \theta e_3 + f(x,t),\ (x,t)\in \Omega \times \mathbb{R}_{\tau}, \\
\partial_{t}\theta+\kappa(-\Delta)^{\beta}\theta+(u\cdot\nabla)\theta = g(x,t),\ (x,t)\in \Omega \times \mathbb{R}_{\tau}, \\
\nabla\cdot u = 0,\ (x,t)\in \Omega \times \mathbb{R}_{\tau},  \\
u(x,\tau) = u_\tau(x), \quad \theta(x,\tau) = \theta_\tau(x),\ x \in \Omega, 
\end{cases}
\end{equation}

\begin{equation}\label{ass:forcing}
G_f := \sup_{r\in\mathbb{R}}
\int_{r-1}^{r}\|f(s)\|_{L^2(\Omega)}^2\,ds < \infty, \quad
G_g := \sup_{r\in\mathbb{R}}
\int_{r-1}^{r}\|g(s)\|_{L^2(\Omega)}^2\,ds < \infty.
\end{equation}

\begin{lemma}\label{lem0301}
Assume that \cref{ass:forcing} holds. For any $\tau \in \mathbb{R}$, there exists a positive constant $\rho_1$ such that for every bounded subset $B_{\ell} \in \mathcal{D}_{\ell}$, there exists a time $\tau_1 =\tau_1(B_{\ell}) \geq 0 $, for all strong solutions to \cref{eq0101} with short trajectory $\chi(s,\tau; (u_\tau, \theta_\tau)) \in B_{\ell}$, we have 
\[
\| u(t) \|_{H_{1}}^{2}
+ \| \theta(t) \|_{H_{2}}^{2}
+ \int_{0}^{\ell} \left(
\| \Lambda^\alpha u(t+\zeta)\|_{L^2(\Omega)}^2 + \| \Lambda^\beta \theta(t+\zeta)\|_{L^2(\Omega)}^2 \right) d\zeta
\leq \rho_1 
\]
for all $t-\tau \geq \tau_1$.
\end{lemma}

\begin{proof}
Multiplying $\eqref{eq0101}_{2}$ by $\theta$ and applying Young's inequality, we get that
\begin{align*}
  \frac{1}{2}\frac{d}{dt}\|\theta(t)\|_{L^2(\Omega)}^2 + \kappa\| \Lambda^\beta \theta(t)\|_{L^2(\Omega)}^2
   & \leq \frac{\kappa\lambda_1^{2\beta}}{2} \|\theta(t)\|_{L^2(\Omega)}^2 + \frac{1}{2\kappa\lambda_1^{2\beta}} \|g(t)\|_{L^2(\Omega)}^2 \\
   & \leq \frac{\kappa}{2} \| \Lambda^\beta \theta(t)\|_{L^2(\Omega)}^2 + \frac{1}{2\kappa\lambda_1^{2\beta}} \|g(t)\|_{L^2(\Omega)}^2,
\end{align*}
since the Poincar\'{e} inequality
$$\| \Lambda^\beta \theta(t)\|_{L^2(\Omega)}^2 \geq \lambda_1^{2\beta}\|\theta(t)\|_{L^2(\Omega)}^2. $$
So we have
\begin{align}\label{eq0301}
  \frac{d}{dt}\|\theta(t)\|_{L^2(\Omega)}^2 + \frac{\kappa\lambda_1^{2\beta}}{2} \|\theta(t)\|_{L^2(\Omega)}^2 + \frac{\kappa}{2}\| \Lambda^\beta \theta(t)\|_{L^2(\Omega)}^2 
   & \leq \frac{1}{\kappa\lambda_1^{2\beta}} \|g(t)\|_{L^2(\Omega)}^2,
\end{align}
Gronwall's inequality subsequently yields that
\begin{align}\label{eq0302} 
  \|\theta(t)\|_{L^2(\Omega)}^2
   & \leq e^{\frac{\kappa\lambda_{1}^{2\beta}}{2}(\tau+\zeta-t)} \|\theta(\tau+\zeta)\|_{L^2(\Omega)}^2 + \frac{1}{\kappa\lambda_{1}^{2\beta}} \int_{\tau+\zeta}^{t} \|g(s)\|_{L^2(\Omega)}^2 e^{\frac{\kappa\lambda_{1}^{2\beta}}{2}(s-t)} ds.
\end{align}
Integrating both sides of \cref{eq0302} with respect to $\zeta$ between $0$ and $\ell$ gives 
\begin{equation}\label{eq0303}
\begin{aligned}
   \|\theta(t)\|_{L^2(\Omega)}^2
   \leq & \frac{1}{\ell} \int_{0}^{\ell} e^{\frac{\kappa\lambda_{1}^{2\beta}}{2}(\tau+\zeta-t)} \|\theta(\tau+\zeta)\|_{L^2(\Omega)}^2 d\zeta \\
   & + \frac{1}{\kappa\lambda_{1}^{2\beta}\ell} \int_{0}^{\ell} \int_{\tau+\zeta}^{t} \|g(s)\|_{L^2(\Omega)}^2 e^{\frac{\kappa\lambda_{1}^{2\beta}}{2}(s-t)} ds d\zeta \\
   \leq & \frac{1}{\ell} e^{\frac{\kappa\lambda_{1}^{2\beta}}{2}(\tau+\ell-t)}  \int_{0}^{\ell} \|\theta(\tau+\zeta)\|_{L^2(\Omega)}^2 d\zeta + \frac{1}{\kappa\lambda_{1}^{2\beta}} \left(1+\frac{2}{\kappa\lambda_{1}^{2\beta}}\right)G_g.
\end{aligned}
\end{equation}
Note that
\begin{align*}
   \int_{\tau+\zeta}^{t} \|g(s)\|_{L^2(\Omega)}^2 e^{\frac{\kappa\lambda_{1}^{2\beta}}{2}(s-t)} ds
   \leq & e^{-\frac{\kappa\lambda_{1}^{2\beta}}{2}t} \sum_{n=0}^{+\infty} \left(\int_{t-(n+1)}^{t-n} \|g(s)\|_{L^2(\Omega)}^2 e^{\frac{\kappa\lambda_{1}^{2\beta}}{2}s} ds\right) \\
   \leq & e^{-\frac{\kappa\lambda_{1}^{2\beta}}{2}t} \sum_{n=0}^{+\infty} e^{\frac{\kappa\lambda_{1}^{2\beta}}{2}(t-n)} \left(\int_{t-(n+1)}^{t-n} \|g(s)\|_{L^2(\Omega)}^2 ds\right) \\
   \leq & \left(1+\frac{2}{\kappa\lambda_{1}^{2\beta}}\right)G_g,
\end{align*}
where
$$
G_g:=\sup_{r \in \mathbb{R}}\int_{r-1}^{r}\|g(s)\|_{L^2(\Omega)}^2 ds< +\infty.
$$
For \cref{eq0303}, there exists a constant $\rho_{1}^{(1)} > 0$ such that for any $B_{\ell} \in \mathcal{D}_{\ell}$, there exists $\tau_{1}^{(1)} = \tau_{1}^{(1)}(B_{\ell}) \geq 0$ with the property that 
\begin{equation}\label{eq0317}
    \|\theta(t)\|_{L^2(\Omega)}^2
   \leq 
   \rho_{1}^{(1)},
\end{equation}
\[   
   \rho_{1}^{(1)}
   :=  
   \frac{2}{\kappa\lambda_{1}^{2\beta}} \left(1+\frac{2}{\kappa\lambda_{1}^{2\beta}}\right)G_g
\]
for all $t \in \mathbb{R}$ with $t - \tau \geq \tau_{1}^{(1)}$. Additionally, integrating \cref{eq0301} both sides between $t$ and $t+\ell$ gives
\begin{equation}\label{eq0307}
\begin{aligned}
    & \frac{\kappa\lambda_1^{2\beta}}{2} \int_{0}^{\ell} \|\theta(t+\zeta)\|_{L^2(\Omega)}^2 d\zeta 
    + 
    \frac{\kappa}{2} \int_{0}^{\ell} \| \Lambda^\beta \theta(t+\zeta)\|_{L^2(\Omega)}^2 d\zeta \\
    \leq & \frac{1}{\kappa\lambda_1^{2\beta}} \int_{t}^{t+\ell} \|g(\zeta)\|_{L^2(\Omega)}^2 d\zeta + \|\theta(t)\|_{L^2(\Omega)}^2 \\
    \leq & \frac{1}{\kappa\lambda_1^{2\beta}} \left(1+[\ell]\right)G_g
    + \frac{1}{\ell} e^{\frac{\kappa\lambda_{1}^{2\beta}}{2}(\tau+\ell-t)} \int_{0}^{\ell} \|\theta(\tau+\zeta)\|_{L^2(\Omega)}^2 d\zeta 
    + \frac{1}{\kappa\lambda_{1}^{2\beta}} \left(1+\frac{2}{\kappa\lambda_{1}^{2\beta}}\right)G_g \\
    \leq & \frac{1}{\ell} e^{\frac{\kappa\lambda_{1}^{2\beta}}{2}(\tau+\ell-t)} \int_{0}^{\ell} \|\theta(\tau+\zeta)\|_{L^2(\Omega)}^2 d\zeta 
    + \frac{1}{\kappa\lambda_{1}^{2\beta}} \left(2+\frac{2}{\kappa\lambda_{1}^{2\beta}}+[\ell]\right)G_g,
\end{aligned}
\end{equation}
where $[\ell]$ denotes the floor function of $\ell$(i.e., the greatest integer less than or equal to $\ell$).
With respect to \cref{eq0307}, there exists a constant $\rho_{1}^{(2)} > 0$ such that for any $B_{\ell} \in \mathcal{D}_{\ell}$, there exists $\tau_{1}^{(2)} = \tau_{1}^{(2)}(B_{\ell}) \geq 0$ ensuring that 
\begin{equation}\label{eq0316}
    \int_{0}^{\ell} \| \Lambda^\beta \theta(t+\zeta)\|_{L^2(\Omega)}^2 d\zeta
    \leq 
    \rho_{1}^{(2)},
\end{equation}
\[
   \rho_{1}^{(2)}
   :=
   \frac{4}{\kappa\lambda_{1}^{2\beta}\min\{\kappa\lambda_1^{2\beta},\kappa\}}
   \left(2+\frac{2}{\kappa\lambda_{1}^{2\beta}}+[\ell]\right)G_g
\]
holds for all $t \in \mathbb{R}$ whenever $t - \tau \geq \tau_{1}^{(2)}$.  

Taking the inner product of $\eqref{eq0101}_{1}$ with $u$ and the Poincar\'{e} inequality
$$\| \Lambda^\alpha u(t)\|_{L^2(\Omega)}^2 \geq \eta_1^{2\alpha}\|u(t)\|_{L^2(\Omega)}^2, $$
we obtain
\begin{align*}
   \frac{1}{2}\frac{d}{dt}\|u(t)\|_{L^2(\Omega)}^2 + \nu\|\Lambda^\alpha u(t)\|_{L^2(\Omega)}^2
   & \leq \frac{\nu\eta_1^{2\alpha}}{2} \|u(t)\|_{L^2(\Omega)}^2 + \frac{1}{\nu\eta_1^{2\alpha}} \|\theta(t)\|_{L^2(\Omega)}^2 + \frac{1}{\nu\eta_1^{2\alpha}} \|f(t)\|_{L^2(\Omega)}^2 \\
   & \leq \frac{\nu}{2} \|\Lambda^\alpha u(t)\|_{L^2(\Omega)}^2 + \frac{1}{\nu\eta_1^{2\alpha}} \left( \|\theta(t)\|_{L^2(\Omega)}^2 +    \|f(t)\|_{L^2(\Omega)}^2 \right).
\end{align*}
Therefore, we get
\begin{align}\label{eq0304}
   \frac{d}{dt}\|u(t)\|_{L^2(\Omega)}^2 + \frac{\nu\eta_1^{2\alpha}}{2} \|u(t)\|_{L^2(\Omega)}^2 + \frac{\nu}{2} \|\Lambda^\alpha u(t)\|_{L^2(\Omega)}^2
   & \leq \frac{2}{\nu\eta_1^{2\alpha}} \left( \|\theta(t)\|_{L^2(\Omega)}^2 +    \|f(t)\|_{L^2(\Omega)}^2 \right).
\end{align}
Using Gronwall's inequality and \cref{eq0303}, we obtain
\begin{align*} 
   \|u(t)\|_{L^2(\Omega)}^2
   \leq 
   & e^{\frac{\nu\eta_1^{2\alpha}}{2}(\tau+\zeta-t)} \|u(\tau+\zeta)\|_{L^2(\Omega)}^2 
   + \frac{2}{\nu\eta_1^{2\alpha}} \int_{\tau+\zeta}^{t} \|\theta(s)\|_{L^2(\Omega)}^2 e^{\frac{\nu\eta_1^{2\alpha}}{2}(s-t)} ds \\
   &+ \frac{2}{\nu\eta_1^{2\alpha}} \int_{\tau+\zeta}^{t}  \|f(s)\|_{L^2(\Omega)}^2 e^{\frac{\nu\eta_1^{2\alpha}}{2}(s-t)} ds \\
   \leq 
   & e^{\frac{\nu\eta_1^{2\alpha}}{2}(\tau+\zeta-t)} \|u(\tau+\zeta)\|_{L^2(\Omega)}^2 
   + \frac{2}{\nu\eta_1^{2\alpha}} \int_{\tau+\zeta}^{t}  \|f(s)\|_{L^2(\Omega)}^2 e^{\frac{\nu\eta_1^{2\alpha}}{2}(s-t)} ds \\
   &
   + \frac{2}{\nu\eta_1^{2\alpha}} \|\theta(\tau+\zeta)\|_{L^2(\Omega)}^2  \int_{\tau+\zeta}^{t} e^{\frac{\kappa\lambda_{1}^{2\beta}}{2}(\tau+\zeta-s)} e^{\frac{\nu\eta_1^{2\alpha}}{2}(s-t)} ds \\
   &
   + \frac{2}{\nu\eta_1^{2\alpha}\kappa\lambda_{1}^{2\beta}} \int_{\tau+\zeta}^{t} \left( \int_{\tau+\zeta}^{s} \|g(\sigma)\|_{L^2(\Omega)}^2 e^{\frac{\kappa\lambda_{1}^{2\beta}}{2}(\sigma-s)} d\sigma \right) e^{\frac{\nu\eta_1^{2\alpha}}{2}(s-t)} ds.
\end{align*}
Note that
\begin{align*}
   \int_{\tau+\zeta}^{t}  \|f(s)\|_{L^2(\Omega)}^2 e^{\frac{\nu\eta_1^{2\alpha}}{2}(s-t)} ds
   \leq \left(1+\frac{2}{\nu\eta_1^{2\alpha}}\right)G_f,
\end{align*}
where
$$G_f:=\sup_{r \in \mathbb{R}}\int_{r-1}^{r}\|f(s)\|_{L^2(\Omega)}^2 ds< +\infty.$$

\noindent{\bfseries Case 1}: $\nu\eta_1^{2\alpha} \neq \kappa\lambda_1^{2\beta}$. In addition, we get
\[
\int_{\tau+\zeta}^{t} e^{\frac{\kappa\lambda_{1}^{2\beta}}{2}(\tau+\zeta-s)} e^{\frac{\nu\eta_1^{2\alpha}}{2}(s-t)} ds = \frac{2}{\nu\eta_1^{2\alpha} - \kappa\lambda_1^{2\beta}} \left( e^{\frac{\kappa\lambda_1^{2\beta}(\tau + \zeta - t)}{2}} - e^{\frac{\nu\eta_1^{2\alpha}(\tau + \zeta - t)}{2}} \right).
\]
Hence, when  \(\nu\eta_1^{2\alpha} \neq \kappa\lambda_1^{2\beta}\) is satisfied, we get
\begin{equation}\label{eq0305}
\begin{aligned} 
   \|u(t)\|_{L^2(\Omega)}^2
   \leq 
   & e^{\frac{\nu\eta_1^{2\alpha}}{2}(\tau+\zeta-t)} \|u(\tau+\zeta)\|_{L^2(\Omega)}^2 
    \\
   &
   + \frac{4}{\nu\eta_1^{2\alpha}(\nu\eta_1^{2\alpha} - \kappa\lambda_1^{2\beta})} \|\theta(\tau+\zeta)\|_{L^2(\Omega)}^2 
   \left( e^{\frac{\kappa\lambda_1^{2\beta}(\tau + \zeta - t)}{2}} - e^{\frac{\nu\eta_1^{2\alpha}(\tau + \zeta - t)}{2}} \right) \\
   &
   + \frac{4}{\nu^{2}\eta_1^{4\alpha}\kappa\lambda_{1}^{2\beta}}
   \left(1+\frac{2}{\kappa\lambda_{1}^{2\beta}}\right)
   G_g
   + \frac{2}{\nu\eta_1^{2\alpha}} \left(1+\frac{2}{\nu\eta_1^{2\alpha}}\right)G_f;
\end{aligned}
\end{equation}
Integrating both sides above with respect to $\zeta$ from $0$ to $\ell$, we can derive that
\begin{equation}\label{eq0308}
\begin{aligned} 
   \|u(t)\|_{L^2(\Omega)}^2
   \leq 
   & \frac{1}{\ell} e^{\frac{\nu\eta_1^{2\alpha}}{2}(\tau+\ell-t)}  \int_{0}^{\ell} \|u(\tau+\zeta)\|_{L^2(\Omega)}^2 d\zeta
    \\
   &
   + \frac{4}{\nu\eta_1^{2\alpha}|\nu\eta_1^{2\alpha} - \kappa\lambda_1^{2\beta}|}  
   \frac{1}{\ell}
   e^{\frac{\min\{\nu\eta_1^{2\alpha},\kappa\lambda_1^{2\beta}\}}{2}(\tau+\ell-t)}
   \int_{0}^{\ell} 
   \|\theta(\tau+\zeta)\|_{L^2(\Omega)}^2  d\zeta \\
   &
   +\frac{4}{\nu^{2}\eta_1^{4\alpha}\kappa\lambda_{1}^{2\beta}}
   \left(1+\frac{2}{\kappa\lambda_{1}^{2\beta}}\right)
   G_g
   + \frac{2}{\nu\eta_1^{2\alpha}} \left(1+\frac{2}{\nu\eta_1^{2\alpha}}\right)G_f.
\end{aligned}
\end{equation}
For \cref{eq0308}, there exists a constant $\rho_{1}^{(3)} > 0$ such that for any $B_{\ell} \in \mathcal{D}_{\ell}$, there is a $\tau_{1}^{(3)} = \tau_{1}^{(3)}(B_{\ell}) \geq 0$ where  
\begin{equation}\label{eq0315}
    \|u(t)\|_{L^2(\Omega)}^2 
    \leq 
    \rho_{1}^{(3)},
\end{equation}
\[
    \rho_{1}^{(3)} 
    := 
    \frac{12}{\nu^{2}\eta_1^{4\alpha}\kappa\lambda_{1}^{2\beta}}
    \left(1+\frac{2}{\kappa\lambda_{1}^{2\beta}}\right)
    G_g
    + 
    \frac{6}{\nu\eta_1^{2\alpha}} \left(1+\frac{2}{\nu\eta_1^{2\alpha}}\right)G_f
\]  
for all $t \in \mathbb{R}$ satisfying $t - \tau \geq \tau_{1}^{(3)}$. By integrating both sides of $\eqref{eq0304}$ from $t$ to $t+\ell$, we have
\begin{equation}\label{eq0309}
\begin{aligned}
  &\frac{\nu\eta_1^{2\alpha}}{2} \int_{0}^{\ell} \|u(t+\zeta)\|_{L^2(\Omega)}^2 d\zeta + \frac{\nu}{2} \int_{0}^{\ell} \| \Lambda^\alpha u(t+\zeta)\|_{L^2(\Omega)}^2 d\zeta \\
  \leq 
  &
  \frac{2}{\nu\eta_1^{2\alpha}} \int_{t}^{t+\ell} \|\theta(\zeta)\|_{L^2(\Omega)}^2 d\zeta 
  +\frac{2}{\nu\eta_1^{2\alpha}} \int_{t}^{t+\ell} \|f(\zeta)\|_{L^2(\Omega)}^2 d\zeta 
  + \|u(t)\|_{L^2(\Omega)}^2 \\
  \leq 
  & 
  \frac{4}{\nu\eta_1^{2\alpha}\kappa\lambda_1^{2\beta}} \frac{1}{\ell} e^{\frac{\kappa\lambda_{1}^{2\beta}}{2}(\tau+\ell-t)} \int_{0}^{\ell} \|\theta(\tau+\zeta)\|_{L^2(\Omega)}^2 d\zeta
  + \frac{1}{\ell} e^{\frac{\nu\eta_1^{2\alpha}}{2}(\tau+\ell-t)}  \int_{0}^{\ell} \|u(\tau+\zeta)\|_{L^2(\Omega)}^2 d\zeta \\
  &
   + \frac{4}{\nu\eta_1^{2\alpha}|\nu\eta_1^{2\alpha} - \kappa\lambda_1^{2\beta}|}  
   \frac{1}{\ell}
   e^{\frac{\min\{\nu\eta_1^{2\alpha},\kappa\lambda_1^{2\beta}\}}{2}(\tau+\ell-t)}
   \int_{0}^{\ell} 
   \|\theta(\tau+\zeta)\|_{L^2(\Omega)}^2  d\zeta \\
  &
  + \frac{2}{\nu\eta_1^{2\alpha}\kappa\lambda_{1}^{2\beta}} 
  \left(1+\frac{2}{\kappa\lambda_{1}^{2\beta}}\right)
  \left(\ell+\frac{2}{\nu\eta_1^{2\alpha}}\right)
  G_g
  + \frac{2}{\nu\eta_1^{2\alpha}} \left(2+[\ell]+\frac{2}{\nu\eta_1^{2\alpha}}\right)G_f.
\end{aligned}
\end{equation}
For \cref{eq0309}, there exists a constant $\rho_{1}^{(4)} > 0$ such that for any $B_{\ell} \in \mathcal{D}_{\ell}$, there exists $\tau_{1}^{(4)} = \tau_{1}^{(4)}(B_{\ell}) \geq 0$ where   
\begin{equation}\label{eq0314}
    \int_{0}^{\ell} \| \Lambda^\alpha u(t+\zeta)\|_{L^2(\Omega)}^2 d\zeta
    \leq 
    \rho_{1}^{(4)},
\end{equation} 
with
\[
    \rho_{1}^{(4)} 
    := 
    \frac{16}{\nu\eta_1^{2\alpha}\min\{\nu\eta_1^{2\alpha},\nu\}}
    \left(
    \frac{1}{\kappa\lambda_{1}^{2\beta}} 
    \left(1+\frac{2}{\kappa\lambda_{1}^{2\beta}}\right)
    \left(\ell+\frac{2}{\nu\eta_1^{2\alpha}}\right)
    G_g
    + 
    \left(2+[\ell]+\frac{2}{\nu\eta_1^{2\alpha}}\right)G_f\right)
\]
for all $t \in \mathbb{R}$ satisfying $t - \tau \geq \tau_{1}^{(4)}$.

\noindent{\bfseries Case 2}: $\nu\eta_1^{2\alpha} = \kappa\lambda_1^{2\beta}$. By a direct calculation, we obtain
\[
\int_{\tau+\zeta}^{t} e^{\frac{\kappa\lambda_{1}^{2\beta}}{2}(\tau+\zeta-s)} e^{\frac{\nu\eta_1^{2\alpha}}{2}(s-t)} ds = (t - \tau - \zeta)e^{\frac{\nu\eta_1^{2\alpha}}{2}(\tau+\zeta - t)}.
\]
When \(\nu\eta_1^{2\alpha} = \kappa\lambda_1^{2\beta}\) holds, the result follows as 
\begin{equation}\label{eq0306}
\begin{aligned} 
   \|u(t)\|_{L^2(\Omega)}^2
   \leq 
   & e^{\frac{\nu\eta_1^{2\alpha}}{2}(\tau+\zeta-t)} \|u(\tau+\zeta)\|_{L^2(\Omega)}^2 
   \\
   &
   + \frac{2}{\nu\eta_1^{2\alpha}} \|\theta(\tau+\zeta)\|_{L^2(\Omega)}^2  (t - \tau - \zeta)e^{\frac{\nu\eta_1^{2\alpha}}{2}(\tau+\zeta - t)} \\
   &
   + \frac{4}{\nu^{3}\eta_1^{6\alpha}} 
   \left(1+\frac{2}{\nu\eta_1^{2\alpha}}\right)G_g
   + \frac{2}{\nu\eta_1^{2\alpha}} \left(1+\frac{2}{\nu\eta_1^{2\alpha}}\right)G_f.
\end{aligned}
\end{equation}
Integrating both sides above with respect to $\zeta$ from $0$ to $\ell$, we get
\begin{equation}\label{eq0310}
\begin{aligned} 
   \|u(t)\|_{L^2(\Omega)}^2
   \leq 
   & \frac{1}{\ell} e^{\frac{\nu\eta_1^{2\alpha}}{2}(\tau+\ell-t)}  \int_{0}^{\ell} \|u(\tau+\zeta)\|_{L^2(\Omega)}^2 d\zeta
    \\
   &
   + \frac{2}{\nu\eta_1^{2\alpha}}
   \frac{1}{\ell}(t-\tau)
   e^{\frac{\nu\eta_1^{2\alpha}}{2}(\tau+\ell-t)}
   \int_{0}^{\ell} 
   \|\theta(\tau+\zeta)\|_{L^2(\Omega)}^2  d\zeta \\
   &
   + \frac{4}{\nu^{3}\eta_1^{6\alpha}} 
   \left(1+\frac{2}{\nu\eta_1^{2\alpha}}\right)
   G_g
   + \frac{2}{\nu\eta_1^{2\alpha}} \left(1+\frac{2}{\nu\eta_1^{2\alpha}}\right)G_f.
\end{aligned}
\end{equation}
Since for all $ t \in \mathbb{R} $, 
\[
(t-\tau)e^{\frac{\nu\eta_1^{2\alpha}}{2}(\tau-t)}
\to 0, \quad\text{as}\quad t - \tau \to \infty, 
\]
it follows that for \cref{eq0310}, there exists a constant $\rho_{1}^{(5)} > 0$ such that for any $B_{\ell} \in \mathcal{D}_{\ell}$, there exists $\tau_{1}^{(5)} = \tau_{1}^{(5)}(B_{\ell}) \geq 0$ such that 
\begin{equation}\label{eq0313}
    \|u(t)\|_{L^2(\Omega)}^2 
    \leq 
    \rho_{1}^{(5)},
\end{equation} 
with
\[
    \rho_{1}^{(5)} 
    := \frac{12}{\nu^{3}\eta_1^{6\alpha}} 
   \left(1+\frac{2}{\nu\eta_1^{2\alpha}}\right)
   G_g
   + \frac{6}{\nu\eta_1^{2\alpha}} \left(1+\frac{2}{\nu\eta_1^{2\alpha}}\right)G_f
\]
for all $t \in \mathbb{R}$ satisfying $t - \tau \geq \tau_{1}^{(5)}$. Integrating $\eqref{eq0304}$ from $t$ to $t+\ell$, we get
\begin{equation}\label{eq0311}
\begin{aligned}
  &\frac{\nu\eta_1^{2\alpha}}{2} \int_{0}^{\ell} \|u(t+\zeta)\|_{L^2(\Omega)}^2 d\zeta + \frac{\nu}{2} \int_{0}^{\ell} \| \Lambda^\alpha u(t+\zeta)\|_{L^2(\Omega)}^2 d\zeta \\
  \leq 
  & 
  \frac{4}{\nu^{2}\eta_1^{4\alpha}} \frac{1}{\ell} e^{\frac{\nu\eta_1^{2\alpha}}{2}(\tau+\ell-t)} \int_{0}^{\ell} \|\theta(\tau+\zeta)\|_{L^2(\Omega)}^2 d\zeta
  + \frac{1}{\ell} e^{\frac{\nu\eta_1^{2\alpha}}{2}(\tau+\ell-t)}  \int_{0}^{\ell} \|u(\tau+\zeta)\|_{L^2(\Omega)}^2 d\zeta \\
  &
  + \frac{2}{\nu\eta_1^{2\alpha}}
   \frac{1}{\ell}(t-\tau)
   e^{\frac{\nu\eta_1^{2\alpha}}{2}(\tau+\ell-t)}
   \int_{0}^{\ell} 
   \|\theta(\tau+\zeta)\|_{L^2(\Omega)}^2  d\zeta \\
  &
  + \frac{2}{\nu^{2}\eta_1^{4\alpha}} 
  \left(1+\frac{2}{\nu\eta_1^{2\alpha}}\right)
  \left(\ell+\frac{2}{\nu\eta_1^{2\alpha}}\right)
  G_g
  + \frac{2}{\nu\eta_1^{2\alpha}} \left(2+[\ell]+\frac{2}{\nu\eta_1^{2\alpha}}\right)G_f.
\end{aligned}
\end{equation}
For \cref{eq0311}, there exists a constant $\rho_{1}^{(6)} > 0$ such that for any $B_{\ell} \in \mathcal{D}_{\ell}$, there exists $\tau_{1}^{(6)} = \tau_{1}^{(6)}(B_{\ell}) \geq 0$ where  
\begin{equation}\label{eq0312}
    \int_{0}^{\ell} \| \Lambda^\alpha u(t+\zeta)\|_{L^2(\Omega)}^2 d\zeta
    \leq 
    \rho_{1}^{(6)},
\end{equation}
with
\[
    \rho_{1}^{(6)} 
    := 
    \frac{16}{\nu\eta_1^{2\alpha}\min\{\nu\eta_1^{2\alpha},\nu\}}
    \left(
    \frac{1}{\nu\eta_1^{2\alpha}} 
    \left(1+\frac{2}{\nu\eta_1^{2\alpha}}\right)
    \left(\ell+\frac{2}{\nu\eta_1^{2\alpha}}\right)
    G_g
    + \left(2+[\ell]+\frac{2}{\nu\eta_1^{2\alpha}}\right)G_f
    \right)
\]
for all $t \in \mathbb{R}$ satisfying $t - \tau \geq \tau_{1}^{(6)}$.
Consequently, for any $\tau \in \mathbb{R}$, there exists a constant 
$$
\rho_1=\rho_{1}^{(1)}+\rho_{1}^{(2)}+\rho_{1}^{(3)}+
\rho_{1}^{(4)}+\rho_{1}^{(5)}+\rho_{1}^{(6)}>0
$$
such that for every $B_{\ell} \in \mathcal{D}_{\ell}$, there exists a time 
$$
\tau_1 =\tau_1(B_{\ell})=\max\{\tau_{1}^{(1)},
            \tau_{1}^{(2)},
            \tau_{1}^{(3)},
            \tau_{1}^{(4)},
            \tau_{1}^{(5)},
            \tau_{1}^{(6)}\}\geq 0 ,
$$
for all strong solutions to \cref{eq0101} with short trajectory $\chi(s,\tau; (u_\tau, \theta_\tau)) \in B_{\ell}$, we obtain
\[
\| u(t) \|_{H_{1}}^{2}
+ \| \theta(t) \|_{H_{2}}^{2}
+ \int_{0}^{\ell} \left(
\| \Lambda^\alpha u(t+\zeta)\|_{L^2(\Omega)}^2 + \| \Lambda^\beta \theta(t+\zeta)\|_{L^2(\Omega)}^2 \right) d\zeta
\leq \rho_1 
\]%
for all $t-\tau \geq \tau_1$. This completes the proof of \cref{lem0301}.
\end{proof}

\end{document}
